
\documentclass{article}
\usepackage{colortbl}
\usepackage{makecell}
\usepackage{multirow}
\usepackage{supertabular}

\begin{document}

\newcounter{utterance}

\twocolumn

{ \footnotesize  \setcounter{utterance}{1}
\setlength{\tabcolsep}{0pt}
\begin{supertabular}{c@{$\;$}|p{.15\linewidth}@{}p{.15\linewidth}p{.15\linewidth}p{.15\linewidth}p{.15\linewidth}p{.15\linewidth}}

    \# & $\;$A & \multicolumn{4}{c}{Game Master} & $\;\:$B\\
    \hline 

    \theutterance \stepcounter{utterance}  

    & & \multicolumn{4}{p{0.6\linewidth}}{\cellcolor[rgb]{0.9,0.9,0.9}{%
	\makecell[{{p{\linewidth}}}]{% 
	  \tt {\tiny [A$\langle$GM]}  
	 State giocando a un gioco di negoziazione cooperativa in cui dovete accordarvi su come dividere una serie di oggetti tra voi e un altro giocatore.\\ \tt \\ \tt Regole:\\ \tt (a) A Lei e all'altro giocatore viene dato un insieme di oggetti. Ognuno di voi riceve anche una funzione di valore segreta, che rappresenta il valore di ciascun tipo di oggetto.\\ \tt (b) Si scambiano messaggi con l'altro giocatore per concordare chi si aggiudica gli oggetti. Potete inviare un massimo di 5 messaggi ciascuno, oppure terminare in anticipo facendo una proposta segreta in qualsiasi momento.\\ \tt (c) A ciascuno di voi viene chiesto di inviare una proposta segreta indicando gli oggetti che desiderate, formattata tra parentesi quadre come segue: "[Proposta: <numero> <nome oggetto>, <numero> <nome oggetto>, <...>]".\\ \tt (d) Se le vostre proposte sono complementari, cioè ci sono abbastanza oggetti per soddisfare entrambe le proposte, a ciascun giocatore viene assegnato un punteggio basato sulla somma dei valori degli oggetti ricevuti. In caso contrario, entrambi ricevono zero punti.\\ \tt (e) L'obiettivo è massimizzare la somma del proprio punteggio e di quello dell'altro giocatore.\\ \tt \\ \tt Cominciamo.\\ \tt \\ \tt L'insieme degli oggetti disponibili è:\\ \tt \\ \tt 1 ditale, 1 sapone, 4 francobolli\\ \tt \\ \tt La funzione valore segreta è:\\ \tt \\ \tt ditale: 3, sapone: 7, francobollo: 0\\ \tt \\ \tt IMPORTANTE: i vostri messaggi, a meno che non si tratti di una proposta segreta, vengono trasmessi direttamente all'altro giocatore, quindi non includete alcuna risposta alle regole o testo di annuncio del vostro messaggio. Per fare una proposta segreta, utilizzate il formato indicato. Non utilizzare delle parentesi quadre quando si comunica all'altro giocatore, altrimenti verrà interpretata come una proposta segreta.\\ \tt  
	  } 
	   } 
	   } 
	 & \\ 
 

    \theutterance \stepcounter{utterance}  

    & \multicolumn{4}{p{0.6\linewidth}}{\cellcolor[rgb]{0.8,1,0.9}{%
	\makecell[{{p{\linewidth}}}]{% 
	  \tt {\tiny [A$\rangle$GM]}  
	 Ciao! Vedo che abbiamo a disposizione 1 ditale, 1 sapone e 4 francobolli da dividere tra noi. Per trovare la divisione migliore, dovremmo condividere le nostre preferenze. \\ \tt \\ \tt Personalmente, il sapone ha molto valore per me, seguito dal ditale. I francobolli invece non mi interessano particolarmente. E tu? Quali oggetti preferisci tra questi? 
	  } 
	   } 
	   } 
	 & & \\ 
 

    \theutterance \stepcounter{utterance}  

    & & & \multicolumn{2}{p{0.3\linewidth}}{\cellcolor[rgb]{0.95,0.95,0.95}{%
	\makecell[{{p{\linewidth}}}]{% 
	  \tt {\tiny [GM$|$GM]}  
	 continue 
	  } 
	   } 
	   } 
	 & & \\ 
 

    \theutterance \stepcounter{utterance}  

    & & & \multicolumn{2}{p{0.3\linewidth}}{\cellcolor[rgb]{0.95,0.95,0.95}{%
	\makecell[{{p{\linewidth}}}]{% 
	  \tt {\tiny [GM$|$GM]}  
	 Ciao! Vedo che abbiamo a disposizione 1 ditale, 1 sapone e 4 francobolli da dividere tra noi. Per trovare la divisione migliore, dovremmo condividere le nostre preferenze. \\ \tt \\ \tt Personalmente, il sapone ha molto valore per me, seguito dal ditale. I francobolli invece non mi interessano particolarmente. E tu? Quali oggetti preferisci tra questi? 
	  } 
	   } 
	   } 
	 & & \\ 
 

    \theutterance \stepcounter{utterance}  

    & & \multicolumn{4}{p{0.6\linewidth}}{\cellcolor[rgb]{0.9,0.9,0.9}{%
	\makecell[{{p{\linewidth}}}]{% 
	  \tt {\tiny [GM$\rangle$B]}  
	 State giocando a un gioco di negoziazione cooperativa in cui dovete accordarvi su come dividere una serie di oggetti tra voi e un altro giocatore.\\ \tt \\ \tt Regole:\\ \tt (a) A Lei e all'altro giocatore viene dato un insieme di oggetti. Ognuno di voi riceve anche una funzione di valore segreta, che rappresenta il valore di ciascun tipo di oggetto.\\ \tt (b) Si scambiano messaggi con l'altro giocatore per concordare chi si aggiudica gli oggetti. Potete inviare un massimo di 5 messaggi ciascuno, oppure terminare in anticipo facendo una proposta segreta in qualsiasi momento.\\ \tt (c) A ciascuno di voi viene chiesto di inviare una proposta segreta indicando gli oggetti che desiderate, formattata tra parentesi quadre come segue: "[Proposta: <numero> <nome oggetto>, <numero> <nome oggetto>, <...>]".\\ \tt (d) Se le vostre proposte sono complementari, cioè ci sono abbastanza oggetti per soddisfare entrambe le proposte, a ciascun giocatore viene assegnato un punteggio basato sulla somma dei valori degli oggetti ricevuti. In caso contrario, entrambi ricevono zero punti.\\ \tt (e) L'obiettivo è massimizzare la somma del proprio punteggio e di quello dell'altro giocatore.\\ \tt \\ \tt Cominciamo.\\ \tt \\ \tt L'insieme degli oggetti disponibili è:\\ \tt \\ \tt 1 ditale, 1 sapone, 4 francobolli\\ \tt \\ \tt La funzione valore segreta è:\\ \tt \\ \tt ditale: 2, sapone: 0, francobollo: 2\\ \tt \\ \tt IMPORTANTE: i vostri messaggi, a meno che non si tratti di una proposta segreta, vengono trasmessi direttamente all'altro giocatore, quindi non includete alcuna risposta alle regole o testo di annuncio del vostro messaggio. Per fare una proposta segreta, utilizzate il formato indicato. Non utilizzare delle parentesi quadre quando si comunica all'altro giocatore, altrimenti verrà interpretata come una proposta segreta.\\ \tt \\ \tt \\ \tt Ciao! Vedo che abbiamo a disposizione 1 ditale, 1 sapone e 4 francobolli da dividere tra noi. Per trovare la divisione migliore, dovremmo condividere le nostre preferenze. \\ \tt \\ \tt Personalmente, il sapone ha molto valore per me, seguito dal ditale. I francobolli invece non mi interessano particolarmente. E tu? Quali oggetti preferisci tra questi? 
	  } 
	   } 
	   } 
	 & \\ 
 

    \theutterance \stepcounter{utterance}  

    & & & \multicolumn{4}{p{0.6\linewidth}}{\cellcolor[rgb]{1,0.85,0.72}{%
	\makecell[{{p{\linewidth}}}]{% 
	  \tt {\tiny [GM$\langle$B]}  
	 Ciao! Interessante, le nostre preferenze sembrano abbastanza diverse. Per me il ditale e i francobolli hanno più valore, mentre il sapone non mi serve molto. \\ \tt \\ \tt Dato che tu preferisci il sapone e io non ne ho bisogno, potresti prenderlo tu. Per quanto riguarda il ditale e i francobolli, visto che entrambi li valutiamo, dovremmo trovare un compromesso. Che ne dici se tu prendi il sapone e una parte dei francobolli, mentre io prendo il ditale e il resto dei francobolli?\\ \tt \\ \tt Quanti francobolli pensi sarebbero giusti per ciascuno di noi? 
	  } 
	   } 
	   } 
	  \\ 
 

    \theutterance \stepcounter{utterance}  

    & & & \multicolumn{2}{p{0.3\linewidth}}{\cellcolor[rgb]{0.95,0.95,0.95}{%
	\makecell[{{p{\linewidth}}}]{% 
	  \tt {\tiny [GM$|$GM]}  
	 continue 
	  } 
	   } 
	   } 
	 & & \\ 
 

    \theutterance \stepcounter{utterance}  

    & & & \multicolumn{2}{p{0.3\linewidth}}{\cellcolor[rgb]{0.95,0.95,0.95}{%
	\makecell[{{p{\linewidth}}}]{% 
	  \tt {\tiny [GM$|$GM]}  
	 Ciao! Interessante, le nostre preferenze sembrano abbastanza diverse. Per me il ditale e i francobolli hanno più valore, mentre il sapone non mi serve molto. \\ \tt \\ \tt Dato che tu preferisci il sapone e io non ne ho bisogno, potresti prenderlo tu. Per quanto riguarda il ditale e i francobolli, visto che entrambi li valutiamo, dovremmo trovare un compromesso. Che ne dici se tu prendi il sapone e una parte dei francobolli, mentre io prendo il ditale e il resto dei francobolli?\\ \tt \\ \tt Quanti francobolli pensi sarebbero giusti per ciascuno di noi? 
	  } 
	   } 
	   } 
	 & & \\ 
 

    \theutterance \stepcounter{utterance}  

    & & \multicolumn{4}{p{0.6\linewidth}}{\cellcolor[rgb]{0.9,0.9,0.9}{%
	\makecell[{{p{\linewidth}}}]{% 
	  \tt {\tiny [A$\langle$GM]}  
	 Ciao! Interessante, le nostre preferenze sembrano abbastanza diverse. Per me il ditale e i francobolli hanno più valore, mentre il sapone non mi serve molto. \\ \tt \\ \tt Dato che tu preferisci il sapone e io non ne ho bisogno, potresti prenderlo tu. Per quanto riguarda il ditale e i francobolli, visto che entrambi li valutiamo, dovremmo trovare un compromesso. Che ne dici se tu prendi il sapone e una parte dei francobolli, mentre io prendo il ditale e il resto dei francobolli?\\ \tt \\ \tt Quanti francobolli pensi sarebbero giusti per ciascuno di noi? 
	  } 
	   } 
	   } 
	 & \\ 
 

    \theutterance \stepcounter{utterance}  

    & \multicolumn{4}{p{0.6\linewidth}}{\cellcolor[rgb]{0.8,1,0.9}{%
	\makecell[{{p{\linewidth}}}]{% 
	  \tt {\tiny [A$\rangle$GM]}  
	 Perfetto! Mi sembra un'ottima base di partenza. Dato che tu dai più valore sia al ditale che ai francobolli, mentre io preferisco principalmente il sapone, potremmo fare così:\\ \tt \\ \tt Tu prendi il ditale (che per te vale di più) e la maggior parte dei francobolli, diciamo 3 francobolli. Io prendo il sapone (che per me è il più importante) e 1 francobollo.\\ \tt \\ \tt In questo modo tu avresti: 1 ditale + 3 francobolli\\ \tt E io avrei: 1 sapone + 1 francobollo\\ \tt \\ \tt Questo ti sembra equo? Oppure preferiresti prendere tutti e 4 i francobolli dato che li valuti molto, e io mi tengo solo il sapone? 
	  } 
	   } 
	   } 
	 & & \\ 
 

    \theutterance \stepcounter{utterance}  

    & & & \multicolumn{2}{p{0.3\linewidth}}{\cellcolor[rgb]{0.95,0.95,0.95}{%
	\makecell[{{p{\linewidth}}}]{% 
	  \tt {\tiny [GM$|$GM]}  
	 continue 
	  } 
	   } 
	   } 
	 & & \\ 
 

    \theutterance \stepcounter{utterance}  

    & & & \multicolumn{2}{p{0.3\linewidth}}{\cellcolor[rgb]{0.95,0.95,0.95}{%
	\makecell[{{p{\linewidth}}}]{% 
	  \tt {\tiny [GM$|$GM]}  
	 Perfetto! Mi sembra un'ottima base di partenza. Dato che tu dai più valore sia al ditale che ai francobolli, mentre io preferisco principalmente il sapone, potremmo fare così:\\ \tt \\ \tt Tu prendi il ditale (che per te vale di più) e la maggior parte dei francobolli, diciamo 3 francobolli. Io prendo il sapone (che per me è il più importante) e 1 francobollo.\\ \tt \\ \tt In questo modo tu avresti: 1 ditale + 3 francobolli\\ \tt E io avrei: 1 sapone + 1 francobollo\\ \tt \\ \tt Questo ti sembra equo? Oppure preferiresti prendere tutti e 4 i francobolli dato che li valuti molto, e io mi tengo solo il sapone? 
	  } 
	   } 
	   } 
	 & & \\ 
 

    \theutterance \stepcounter{utterance}  

    & & \multicolumn{4}{p{0.6\linewidth}}{\cellcolor[rgb]{0.9,0.9,0.9}{%
	\makecell[{{p{\linewidth}}}]{% 
	  \tt {\tiny [GM$\rangle$B]}  
	 Perfetto! Mi sembra un'ottima base di partenza. Dato che tu dai più valore sia al ditale che ai francobolli, mentre io preferisco principalmente il sapone, potremmo fare così:\\ \tt \\ \tt Tu prendi il ditale (che per te vale di più) e la maggior parte dei francobolli, diciamo 3 francobolli. Io prendo il sapone (che per me è il più importante) e 1 francobollo.\\ \tt \\ \tt In questo modo tu avresti: 1 ditale + 3 francobolli\\ \tt E io avrei: 1 sapone + 1 francobollo\\ \tt \\ \tt Questo ti sembra equo? Oppure preferiresti prendere tutti e 4 i francobolli dato che li valuti molto, e io mi tengo solo il sapone? 
	  } 
	   } 
	   } 
	 & \\ 
 

    \theutterance \stepcounter{utterance}  

    & & & \multicolumn{4}{p{0.6\linewidth}}{\cellcolor[rgb]{1,0.85,0.72}{%
	\makecell[{{p{\linewidth}}}]{% 
	  \tt {\tiny [GM$\langle$B]}  
	 La tua proposta mi sembra molto ragionevole! Apprezzo che tu abbia considerato le nostre diverse preferenze.\\ \tt \\ \tt Tra le due opzioni che hai suggerito, preferirei la prima: io prendo 1 ditale e 3 francobolli, tu prendi 1 sapone e 1 francobollo. Mi sembra una divisione più equilibrata rispetto a prendere tutti i francobolli, così entrambi otteniamo qualcosa che valutiamo.\\ \tt \\ \tt Se per te va bene questa divisione, sono pronto a formalizzarla. Tu che dici? 
	  } 
	   } 
	   } 
	  \\ 
 

    \theutterance \stepcounter{utterance}  

    & & & \multicolumn{2}{p{0.3\linewidth}}{\cellcolor[rgb]{0.95,0.95,0.95}{%
	\makecell[{{p{\linewidth}}}]{% 
	  \tt {\tiny [GM$|$GM]}  
	 continue 
	  } 
	   } 
	   } 
	 & & \\ 
 

    \theutterance \stepcounter{utterance}  

    & & & \multicolumn{2}{p{0.3\linewidth}}{\cellcolor[rgb]{0.95,0.95,0.95}{%
	\makecell[{{p{\linewidth}}}]{% 
	  \tt {\tiny [GM$|$GM]}  
	 La tua proposta mi sembra molto ragionevole! Apprezzo che tu abbia considerato le nostre diverse preferenze.\\ \tt \\ \tt Tra le due opzioni che hai suggerito, preferirei la prima: io prendo 1 ditale e 3 francobolli, tu prendi 1 sapone e 1 francobollo. Mi sembra una divisione più equilibrata rispetto a prendere tutti i francobolli, così entrambi otteniamo qualcosa che valutiamo.\\ \tt \\ \tt Se per te va bene questa divisione, sono pronto a formalizzarla. Tu che dici? 
	  } 
	   } 
	   } 
	 & & \\ 
 

    \theutterance \stepcounter{utterance}  

    & & \multicolumn{4}{p{0.6\linewidth}}{\cellcolor[rgb]{0.9,0.9,0.9}{%
	\makecell[{{p{\linewidth}}}]{% 
	  \tt {\tiny [A$\langle$GM]}  
	 La tua proposta mi sembra molto ragionevole! Apprezzo che tu abbia considerato le nostre diverse preferenze.\\ \tt \\ \tt Tra le due opzioni che hai suggerito, preferirei la prima: io prendo 1 ditale e 3 francobolli, tu prendi 1 sapone e 1 francobollo. Mi sembra una divisione più equilibrata rispetto a prendere tutti i francobolli, così entrambi otteniamo qualcosa che valutiamo.\\ \tt \\ \tt Se per te va bene questa divisione, sono pronto a formalizzarla. Tu che dici? 
	  } 
	   } 
	   } 
	 & \\ 
 

    \theutterance \stepcounter{utterance}  

    & \multicolumn{4}{p{0.6\linewidth}}{\cellcolor[rgb]{0.8,1,0.9}{%
	\makecell[{{p{\linewidth}}}]{% 
	  \tt {\tiny [A$\rangle$GM]}  
	 Perfetto! Sono completamente d'accordo con questa divisione. Mi sembra davvero equa e rispetta le preferenze di entrambi.\\ \tt \\ \tt Quindi confermiamo:\\ \tt - Tu: 1 ditale + 3 francobolli\\ \tt - Io: 1 sapone + 1 francobollo\\ \tt \\ \tt Procediamo con la formalizzazione di questo accordo!\\ \tt \\ \tt [Proposta: 1 sapone, 1 francobollo] 
	  } 
	   } 
	   } 
	 & & \\ 
 

    \theutterance \stepcounter{utterance}  

    & & & \multicolumn{2}{p{0.3\linewidth}}{\cellcolor[rgb]{0.95,0.95,0.95}{%
	\makecell[{{p{\linewidth}}}]{% 
	  \tt {\tiny [GM$|$GM]}  
	 proposal 
	  } 
	   } 
	   } 
	 & & \\ 
 

    \theutterance \stepcounter{utterance}  

    & & & \multicolumn{2}{p{0.3\linewidth}}{\cellcolor[rgb]{0.95,0.95,0.95}{%
	\makecell[{{p{\linewidth}}}]{% 
	  \tt {\tiny [GM$|$GM]}  
	 [0, 1, 1] 
	  } 
	   } 
	   } 
	 & & \\ 
 

    \theutterance \stepcounter{utterance}  

    & & \multicolumn{4}{p{0.6\linewidth}}{\cellcolor[rgb]{0.9,0.9,0.9}{%
	\makecell[{{p{\linewidth}}}]{% 
	  \tt {\tiny [GM$\rangle$B]}  
	 L'altro giocatore ha presentato una proposta segreta. Ora tocca a Lei presentare una sua proposta.\\ \tt  
	  } 
	   } 
	   } 
	 & \\ 
 

    \theutterance \stepcounter{utterance}  

    & & & \multicolumn{4}{p{0.6\linewidth}}{\cellcolor[rgb]{1,0.85,0.72}{%
	\makecell[{{p{\linewidth}}}]{% 
	  \tt {\tiny [GM$\langle$B]}  
	 [Proposta: 1 ditale, 3 francobolli] 
	  } 
	   } 
	   } 
	  \\ 
 

    \theutterance \stepcounter{utterance}  

    & & & \multicolumn{2}{p{0.3\linewidth}}{\cellcolor[rgb]{0.95,0.95,0.95}{%
	\makecell[{{p{\linewidth}}}]{% 
	  \tt {\tiny [GM$|$GM]}  
	 proposal 
	  } 
	   } 
	   } 
	 & & \\ 
 

    \theutterance \stepcounter{utterance}  

    & & & \multicolumn{2}{p{0.3\linewidth}}{\cellcolor[rgb]{0.95,0.95,0.95}{%
	\makecell[{{p{\linewidth}}}]{% 
	  \tt {\tiny [GM$|$GM]}  
	 [1, 0, 3] 
	  } 
	   } 
	   } 
	 & & \\ 
 

    \theutterance \stepcounter{utterance}  

    & & & \multicolumn{2}{p{0.3\linewidth}}{\cellcolor[rgb]{0.95,0.95,0.95}{%
	\makecell[{{p{\linewidth}}}]{% 
	  \tt {\tiny [GM$|$GM]}  
	 [[0, 1, 1], [1, 0, 3]] 
	  } 
	   } 
	   } 
	 & & \\ 
 

\end{supertabular}
}

\end{document}
